% Učitelská příprava Work/MAT: 1. ročník, hodina 17 podle upraveného rozvržení bloku MNZ.
% Nové pořadí: 12 množina, 13 podmnožiny a operace, 14 diagramy,
% 15 vlastnosti + volitelná mohutnost, 16 intervaly, 17 slovní úlohy.
% Učebnice: ucebnice/01_Zakladni_poznatky/Mnoziny/Množiny.tex
% (Mohutnost množiny, Množinové operace, Vennovy diagramy).
% Vlastní projekt: původní 19_MNZ_SlovniUlohy.tex (obsahová návaznost),
% původní 20_MNZ_OpakovaniMnozin.tex (diagnostika), nová hodina 14.
% Realisticky.cz: 010306 Řešení slovních úloh pomocí Vennových diagramů I;
% 010307 Řešení slovních úloh pomocí Vennových diagramů II.
% Pozn.: přímé PDF lekcí nebylo dostupné (502); použity ověřené veřejné
% výpisy témat a ukázkových zadání, konkrétní úlohy a řešení jsou nové.
\documentclass[11pt,a4paper]{article}
% Spolecny styl pro ucitelske pripravy MAT (Work/MAT)
% Verze 1.1, 2026-10-23
\NeedsTeXFormat{LaTeX2e}
\ProvidesPackage{priprava}[2026/10/23 v1.1 Sablona ucitelskych priprav MAT]

% Zakladni typografie a cestina
\RequirePackage[T1]{fontenc}
\RequirePackage[utf8]{inputenc}
\RequirePackage[czech]{babel}
\RequirePackage{lmodern}
\RequirePackage{microtype}
\RequirePackage{amsmath,amssymb}
\RequirePackage{xcolor}
\RequirePackage{enumitem}
\RequirePackage[most]{tcolorbox}
\RequirePackage{tabularx,booktabs,array}
\RequirePackage{fancyhdr}
\RequirePackage{needspace}
\RequirePackage{geometry}
\RequirePackage[hidelinks]{hyperref}

\geometry{a4paper,top=17mm,bottom=18mm,left=23mm,right=23mm}
\definecolor{notegray}{gray}{0.39}
\setlength{\parindent}{0pt}
\setlength{\parskip}{2.5pt plus 1pt minus 0.5pt}
\setlist[itemize]{topsep=2pt,itemsep=1pt,parsep=0pt,leftmargin=1.45em}
\setlist[enumerate]{topsep=2pt,itemsep=2pt,parsep=0pt,leftmargin=1.65em}
\setcounter{secnumdepth}{0}
\renewcommand{\arraystretch}{1.12}

% Zahlavi: vlevo rocnik | globalni tema, vpravo cislo: nazev hodiny
% Zapati zustava zcela prazdne.
\pagestyle{fancy}
\fancyhf{}
\renewcommand{\headrulewidth}{0.35pt}
\renewcommand{\footrulewidth}{0pt}
\setlength{\headheight}{14pt}
\newcommand{\nastavPripravu}[4]{%
  \fancyhead[L]{\small #1\enspace\textbar\enspace #2}%
  \fancyhead[R]{\small #3: #4}%
}

% Minimalisticke cernobile boxy. Kazda definice zvlast.
\tcbset{commonlesson/.style={enhanced,breakable=false,colback=white,colframe=black,
  colbacktitle=white,coltitle=black,boxrule=0.45pt,
  arc=0.8mm,outer arc=0.8mm,left=3mm,right=3mm,top=1.6mm,bottom=1.6mm,
  toptitle=1mm,bottomtitle=0.7mm,fonttitle=\bfseries,
  attach boxed title to top left={xshift=2.5mm,yshift=-1.15mm},
  boxed title style={colback=white,colframe=white,boxrule=0pt,left=1mm,right=1mm}}}
\newtcolorbox{zadaniBox}[1][]{commonlesson,title={Zadání},#1}
\newtcolorbox{prikladBox}[1][]{commonlesson,title={Příklad},#1}
\newtcolorbox{definiceBox}[1][]{commonlesson,title={Definice},#1}

% Spolecny box pro zadani a jeho reseni, bez mezery mezi castmi.
% Pouziti:
%   \begin{prikladReseniBox}
%     Zadani ulohy.
%     \tcblower
%     \begin{align*} ... \end{align*}
%   \end{prikladReseniBox}
% Spodni cast ma jemny sedy podklad a delici linku; lze v ni sazet
% vice radku matematiky. Obe casti maji jen jeden spolecny ramecek.
\tcbset{solutionlesson/.style={
  commonlesson,
  enhanced jigsaw,
  breakable,
  colbacklower=black!3,
  segmentation style={solid,draw=black!25,line width=0.35pt},
  fontlower=\small\color{notegray},
  before lower={\noindent\textbf{Řešení:}\enspace},
  leftlower=3mm,rightlower=3mm
}}
\newtcolorbox{prikladReseniBox}[1][]{solutionlesson,title={Příklad},#1}
\newtcolorbox{zadaniReseniBox}[1][]{solutionlesson,title={Zadání},#1}

% Pomucky pro autora pripravy: sedive poznamky; \reseni pro starsi pripravy
\newcommand{\poznamka}[1]{\par\begingroup\small\color{notegray}\itshape
  \textnormal{\textbf{Poznámka pro učitele.}}\enspace #1\par\endgroup}
\newcommand{\reseni}[1]{\par\begingroup\small\color{notegray}
  \textbf{Očekávané řešení:}\enspace #1\par\endgroup}
\newcommand{\blok}[1]{\needspace{5\baselineskip}\subsection*{#1}}
\newcommand{\vyrok}[1]{\uv{#1}}

\usepackage{tikz}
\usetikzlibrary{calc,patterns}
\nastavPripravu{1. ročník}{Množiny}{17}{Slovní úlohy: dvě množiny}
\begin{document}

\blok{Od slovní úlohy k počtům v diagramu}
Ve třídě se někteří žáci věnují programování, někteří robotice, někteří
obojímu a další ani jedné činnosti. Jak poznáme, že jsme nikoho
nezapočítali dvakrát? Vennův diagram už známe z minulé hodiny; nyní
\emph{do jednotlivých oblastí nezapisujeme jména, ale počty lidí}.

\poznamka{Základ hodiny: motivace a P1 12 min, odvození vzorce a význam slov 8 min,
P2 8 min, S1 a)--b) s kontrolou 12 min, organizace a závěr 5 min.
Úlohy S2, R1, R2 jsou rozšířením, nikoliv podmínkou zvládnutí hodiny.}
\begin{definiceBox}
\textbf{Mohutnost konečné množiny} $A$, značená $|A|$, je počet jejích
různých prvků. Například $|\{1,2,2,3\}|=3$ a $|\varnothing|=0$.
\end{definiceBox}

V úlohách označuje $U$ všechny sledované osoby nebo předměty.
Množiny $A,B\subseteq U$ označují dvě vlastnosti. Nezaměňuj $A\cap B$
(množina) a $|A\cap B|$ (počet jejích prvků).

\begin{prikladReseniBox}[title={P1: Dva školní kroužky}]
V ročníku je $32$ studentů. Programování navštěvuje $18$ studentů,
robotiku $15$ a oba kroužky $7$ studentů. Urči, kolik studentů navštěvuje:
\textbf{a)} pouze programování; \textbf{b)} pouze robotiku;
\textbf{c)} alespoň jeden kroužek; \textbf{d)} ani jeden.
\tcblower
Označme $A$ programování a $B$ robotiku. Začni společnou částí $7$:
\[
\begin{aligned}
\text{pouze }A &:18-7=11, & \text{pouze }B &:15-7=8,\\
\text{alespoň jeden} &:11+7+8=26, & \text{ani jeden} &:32-26=6.
\end{aligned}
\]
\begin{center}
\begin{tikzpicture}[scale=.88]
\draw[thick] (-2.9,-1.7) rectangle (3,1.85);
\draw[thick] (-.83,0) circle(1.27); \draw[thick] (.83,0) circle(1.27);
\node at (-1.42,1.5) {$A$}; \node at (1.42,1.5) {$B$};
\node[anchor=north east] at (2.88,1.74) {$U$};
\node at (-1.15,0) {\Large $11$}; \node at (0,0) {\Large $7$};
\node at (1.15,0) {\Large $8$}; \node at (2.42,-1.28) {\Large $6$};
\node[font=\footnotesize] at (0,-2.05) {Čísla označují počty v navzájem disjunktních oblastech.};
\end{tikzpicture}
\end{center}
\textbf{Kontrola:} $11+7+8+6=32$, dále $11+7=18$ a $7+8=15$.
Nestačí ověřit jen celkový součet: ověř také počty v obou kroužcích.

\textbf{Čtyřkrokový postup:} (1) vyznač základní množinu $U$,
(2) doplň průnik, (3) dopočti oblasti \uv{jen}, (4) oblast mimo kruhy
a zkontroluj všechny původní údaje.
\end{prikladReseniBox}

\poznamka{Před odhalením řešení požaduj čtyři oddělené oblasti.
Typická chyba: zapsat do části \uv{jen $A$} rovnou $18$.
Ptej se: \uv{Je student z průniku také studentem programování?}}

\blok{Proč společnou část jednou odečítáme?}
Sčítáme-li $|A|+|B|$, započítáme každý společný prvek dvakrát.
Každý prvek ve sjednocení však máme započítat jen jednou.

\begin{definiceBox}
\textbf{Princip inkluze a exkluze pro dvě konečné množiny:}
\[
\boxed{|A\cup B|=|A|+|B|-|A\cap B|.}
\]
\end{definiceBox}

Pro P1 dostáváme $18+15-7=26$. Nejde o další pravidlo, které se musí
mechanicky naučit: vzorec plyne přímo z rozdělení na tři nepřekrývající
se oblasti. Pokud $A\cap B=\varnothing$, nic neodečítáme.

\begin{center}
\begin{tabularx}{\linewidth}{@{}lX@{}}
\toprule
Výraz v zadání & Jaké oblasti hledáme?\\ \midrule
\emph{oba}, \emph{současně} & $A\cap B$\\
\emph{alespoň jeden}, \emph{nebo} & $A\cup B$; společnou část zahrnujeme\\
\emph{právě jeden}, \emph{pouze jeden} & jen $A$ a jen $B$; bez průniku\\
\emph{ani jeden}, \emph{žádný} & $U\setminus(A\cup B)$\\
\emph{nejvýše jeden} & žádný nebo právě jeden; vylučujeme průnik\\
\bottomrule
\end{tabularx}
\end{center}

\begin{prikladReseniBox}[title={P2: Kontrola výrobků -- neznámý překryv}]
Kontrolor má $42$ součástek. Poškrábaný povrch má $25$ kusů,
chybný rozměr $21$ kusů a \emph{ani jednu z těchto vad} nemá $8$ kusů.
Urči počet se dvěma vadami, pouze s jednou vadou a právě s jednou vadou.
\tcblower
Alespoň jednu uvedenou vadu má $42-8=34$ součástek.
Položme $A=$ poškrábané, $B=$ chybný rozměr. Ze vzorce:
\[
|A\cap B|=25+21-34=12.
\]
Jen povrch: $25-12=13$; jen rozměr: $21-12=9$.
Právě jednu vadu má $13+9=22$ kusů. Obě vady má $12$ kusů.
\textbf{Kontrola:} $13+12+9+8=42$.
Údaj \uv{bez uvedených vad} nevylučuje jiné, v zadání nesledované vady.
\end{prikladReseniBox}

\poznamka{Nechej žáky označit všechna známá čísla ještě před počítáním.
U P2 je první krok doplnění \uv{alespoň jedna} z údaje \uv{ani jedna}.
Do povinného jádra patří P1, odvození vzorce a P2; náročnější úlohy jsou níže.}

\blok{Samostatná práce: od údajů k odpovědím}
\begin{zadaniReseniBox}[title={S1: dvě základní úlohy}]
\begin{enumerate}[label=\alph*)]
\item V $36$členné skupině používá CAD $22$ žáků, 3D tisk $19$ žáků a obojí $9$.
Kolik používá pouze CAD, pouze 3D tisk, alespoň jedno a ani jedno?
\item Z $54$ výrobků splňuje alespoň jednu ze dvou vlastností $45$ kusů.
První vlastnost splňuje $28$ kusů a druhou $30$. Kolik má obě vlastnosti?
Kolik nemá žádnou?
\end{enumerate}
\tcblower
\begin{enumerate}[label=\alph*)]
\item Jen CAD $22-9=13$, jen tisk $19-9=10$; alespoň jedno
$13+9+10=32$, ani jedno $36-32=4$. Kontrola $13+9+10+4=36$.
\item Obě $28+30-45=13$, žádnou $54-45=9$.
Jen první $15$, jen druhou $17$; kontrola $15+13+17+9=54$.
\end{enumerate}
\end{zadaniReseniBox}

\begin{zadaniReseniBox}[title={S2: složitější význam slov \uv{právě} a \uv{možné}},breakable=false]
\begin{enumerate}[label=\alph*)]
\item V $48$členné skupině sportuje v oddílu $A$ celkem $30$ žáků,
v oddílu $B$ celkem $26$ a \emph{právě jeden} oddíl navštěvuje $32$ žáků.
Kolik navštěvuje oba oddíly a kolik ani jeden?
\item Ve skupině $30$ lidí má vlastnost $A$ celkem $24$,
vlastnost $B$ celkem $18$ a obě vlastnosti $7$. Jsou údaje možné?
\end{enumerate}
\tcblower
\begin{enumerate}[label=\alph*)]
\item Označme $x$ počet žáků v obou oddílech. Právě jeden:
$(30-x)+(26-x)=32$, tedy $56-2x=32$ a $x=12$.
Sjednocení má $30+26-12=44$ žáků, mimo oba jsou $48-44=4$.
\item Sjednocení by mělo $24+18-7=35$ osob, což je více než $30$.
Údaje jsou v rozporu; nelze je všechny splnit současně.
\end{enumerate}
\end{zadaniReseniBox}

\poznamka{Pro základ vyber S1 a), b). S2 je rozšíření pro rychlejší třídu.
Rozdíl mezi \uv{právě jeden} a \uv{alespoň jeden} není jen slovní;
S2 b) navíc testuje smysluplnost výsledku.}


\blok{Rozšiřující úlohy: dvě množiny a neznámé údaje}
\poznamka{Příklady níže nejsou povinným jádrem. Mohou sloužit k diferenciaci
nebo k dokončení doma. Úlohy se třemi množinami byly přesunuty do 18. hodiny.}
\begin{zadaniReseniBox}[title={R1: když známe počet právě v jedné skupině}]
Ve skupině $48$ osob je $30$ členů klubu $A$, $26$ členů klubu $B$
a právě v jednom klubu je $32$ osob. Kolik je v obou klubech a kolik v žádném?
\tcblower
Označme společnou část $x$. Ve výlučných oblastech je $(30-x)$ a $(26-x)$:
\[
(30-x)+(26-x)=32\quad\Rightarrow\quad 56-2x=32\quad\Rightarrow\quad x=12.
\]
Alespoň jeden klub navštěvuje $30+26-12=44$ osob, žádný $48-44=4$.
Kontrola oblastí: $18+12+14+4=48$.
\end{zadaniReseniBox}
\begin{zadaniReseniBox}[title={R2: ne každé číslo v zadání je možné},breakable=false]
Ve skupině $30$ osob má vlastnost $A$ celkem $24$,
vlastnost $B$ celkem $18$ a obě vlastnosti má $7$ lidí.
Mohou všechny údaje platit současně?
\tcblower
Nemohou. Sjednocení by obsahovalo $24+18-7=35$ osob,
což převyšuje celkový počet $30$. Výsledek není možné vyznačit do diagramu.
\end{zadaniReseniBox}

\blok{Domácí procvičení: pouze dvě množiny}
\begin{zadaniReseniBox}[title={D1: tři situace, vždy čtyři oblasti}]
\begin{enumerate}[label=\alph*)]
\item Z $38$ žáků jezdí do školy na kole $21$, autobusem $18$,
obojí využívá $8$. Urči pouze kolo, pouze autobus, alespoň jednu
možnost a žádnou z uvedených možností.
\item Ve skladu je $60$ výrobků; vadu $A$ má $32$, vadu $B$ má $25$,
ani jednu nemá $15$. Kolik má obě vady a kolik právě jednu?
\item Ve skupině $26$ osob má vlastnost $A$ $16$, vlastnost $B$ $14$
a právě jednu z nich má $20$. Kolik osob má obě a kolik žádnou?
\end{enumerate}
\tcblower
\begin{enumerate}[label=\alph*)]
\item Jen kolo $13$, jen autobus $10$, alespoň jedna $31$, žádná $7$.
Kontrola $13+8+10+7=38$.
\item Alespoň jedna $60-15=45$, obě $32+25-45=12$.
Jen $A$ $20$, jen $B$ $13$, právě jedna $33$.
\item Pro průnik $x$ platí $(16-x)+(14-x)=20$, tedy $x=5$.
Sjednocení $16+14-5=25$, žádná $26-25=1$.
Jen $A$ $11$, jen $B$ $9$; kontrola $11+5+9+1=26$.
\end{enumerate}
\end{zadaniReseniBox}

\end{document}
